\documentclass[11pt,letterpaper]{article}
\usepackage[utf8]{inputenc}
\usepackage[T1]{fontenc}
\usepackage[spanish,es-tabla,es-noshorthands]{babel}
\usepackage{amsmath,amssymb,graphicx,fancyhdr,anysize,float,listings,booktabs,tabularx,xcolor,multirow,caption}
\usepackage{tikz}
\usetikzlibrary{arrows.meta,positioning,fit,calc,babel,backgrounds,decorations.pathreplacing,shapes.geometric}
\usepackage{pgfplots}
\renewcommand{\_}{\textunderscore\allowbreak}
\newcolumntype{Y}{>{\raggedright\arraybackslash}X}
\newcolumntype{P}[1]{>{\raggedright\arraybackslash}p{#1}}
\pgfplotsset{compat=1.17}
\usepackage[hyperindex=true,breaklinks=true,colorlinks=true,linkcolor=black,urlcolor=blue]{hyperref}
\pagestyle{fancy}\fancyhf{}
\chead{Arquitectura de Computadores\vspace{0.3cm}}
\IfFileExists{uandes.pdf}{\lhead{\includegraphics[scale=0.4]{uandes.pdf}}}{}
\lfoot{\scriptsize Facultad de Ingeniería y Ciencias Aplicadas\\Universidad de los Andes}
\rfoot{\thepage}\renewcommand{\footrulewidth}{0.25pt}\addtolength{\headheight}{1.4cm}
\renewcommand{\baselinestretch}{1.25}

% ---------- Colores y estilos (mismos que el Proyecto 3) ----------
\definecolor{azul}{RGB}{0,84,147}
\definecolor{celeste}{RGB}{0,176,216}
\definecolor{verde}{RGB}{46,139,87}
\definecolor{rojo}{RGB}{192,57,43}
\definecolor{naranjo}{RGB}{230,126,34}
\definecolor{gris}{RGB}{110,110,110}
\definecolor{codebg}{RGB}{247,247,250}
\lstdefinestyle{asm}{basicstyle=\ttfamily\footnotesize,keywordstyle=\color{azul}\bfseries,commentstyle=\color{gris}\itshape,numbers=left,numberstyle=\tiny\color{gris},backgroundcolor=\color{codebg},frame=single,rulecolor=\color{gris!40},columns=fullflexible,keepspaces=true,breaklines=true,morecomment=[l]{\#}}
\lstdefinestyle{shell}{basicstyle=\ttfamily\footnotesize,backgroundcolor=\color{codebg},frame=single,rulecolor=\color{gris!40},columns=fullflexible,keepspaces=true,breaklines=true}
\renewcommand{\lstlistingname}{Código}
\lstset{basicstyle=\ttfamily\small,columns=fullflexible,keepspaces=true,breaklines=true,frame=single}
\newcommand{\ok}{\textcolor{verde}{\textbf{Cumple}}}
\newcommand{\nok}{\textcolor{rojo}{\textbf{No cumple}}}
\newcommand{\figopt}[3]{\IfFileExists{#1}{\begin{figure}[H]\centering\includegraphics[width=#2\linewidth]{#1}\caption{#3}\end{figure}}{}}

\begin{document}
\marginsize{2.5cm}{1.5cm}{0.8cm}{2cm}
\begin{titlepage}\begin{center}
\IfFileExists{Log.png}{\includegraphics[width=1.2in]{Log.png}\\}{}
\large UNIVERSIDAD DE LOS ANDES\\
\large Facultad de Ingeniería y Ciencias Aplicadas\\
\large Arquitectura de Computadores 2026-20\\
\vspace{2cm}
{\LARGE \textbf{Proyecto 4\\[0.35cm]{\Large Pochoco Trader: co-diseño hardware/software sobre Espino Core}}}
\vspace{1.5cm}\\
Integrantes:\\ Dariagny Fernández\\ Daniela Rangel\\ Javiera Sepúlveda\\
Profesor:\\ Jorge Gómez
\vfill\today
\end{center}\end{titlepage}

\tableofcontents
\newpage

% =====================================================================
\section{Introducción}
El Proyecto 4 pide diseñar un agente de inversión que se ejecute en el Espino Core sobre la Go Board. Una FPGA maestra transmite por SPI el precio de una acción tick a tick; en cada tick el programa responde, en un byte, si compra, vende o mantiene y cuántas acciones. El patrimonio final en tres rondas, cada una con un período entre precios más corto, define el resultado de la competencia, de modo que la eficiencia del programa y del hardware se transforma en ganancia o pérdida.

El trabajo se organiza en tres bloques:
\begin{enumerate}
  \item \textbf{Algoritmo:} diseño y validación fuera de la placa con la serie de práctica y comparación con dos referencias (no operar y comprar y mantener).
  \item \textbf{Assembly:} implementación en RV32E sobre un Espino sin multiplicación ni shifts, con verificación bit a bit contra la referencia y mediciones por tick.
  \item \textbf{Co-diseño:} tres variantes (software puro, acelerador MMIO y shifter parcial en el ALU), comparadas en LUTs, BRAM, frecuencia máxima, ciclos, CPI, Ops/s y roofline, y justificación de la variante final.
\end{enumerate}

Toda la evidencia de este informe se puede reproducir con los comandos de la Sección~\ref{sec:repro}. Las mediciones provienen de simulación RTL con el mismo protocolo de la maestra. La demostración con la FPGA verifica la integración física; las métricas cuantitativas se obtienen de los testbenches reproducibles del repositorio.

% =====================================================================
\section{Algoritmo de predicción y decisión}

\subsection{Protocolo y consecuencias para el diseño}
Dos propiedades del protocolo condicionan el algoritmo:
\begin{itemize}
  \item \textbf{Retardo de un tick.} El byte que sale en la transacción $n$ es el que estaba en \texttt{DECISION} cuando bajó \texttt{CS\_n}, es decir, la respuesta a $p_{n-1}$; la maestra lo ejecuta al precio $p_n$. Toda señal opera con un tick de retardo y las señales de muy corto plazo no son aprovechables.
  \item \textbf{Cantidad acotada y recorte por la maestra.} Una orden mueve a lo más 63 acciones; una compra se recorta a lo que alcanza la caja y una venta a lo que se tiene, y una orden de cantidad efectiva cero no se ejecuta ni paga comisión.
\end{itemize}

\subsection{Análisis de la serie de práctica}
La serie pública contiene 20\,000 precios normalizados a $[100,255]$. Se usa para validar y comparar implementaciones, pero ninguno de sus valores está incorporado en el programa. Se observaron tres hechos relevantes:
\begin{itemize}
  \item \textbf{El ruido de corto plazo es alto.} Una decisión basada solo en el último cambio de precio produce muchas entradas y salidas y acumula comisiones.
  \item \textbf{Existen movimientos abruptos.} Un salto grande contiene información que una media lenta tarda en reflejar; por eso se trata como una señal separada, sin depender de una posición fija dentro de la serie.
  \item \textbf{Conviene evitar transacciones pequeñas repetidas.} Como cada orden ejecutada paga comisión, la estrategia usa estados persistentes y cambia de posición solo cuando la señal es suficientemente fuerte.
\end{itemize}
Conclusión de diseño: el agente combina una tendencia de mediano plazo con detección de saltos, y solo puede estar \emph{largo} (comprando) o en \emph{efectivo} (vendiendo), pues el protocolo no permite ventas en corto.

\subsection{Estrategia final: tendencia, histéresis y saltos}
La versión final (V2) usa dos medias móviles exponenciales con $\alpha_r=2^{-6}$ (rápida) y $\alpha_l=2^{-8}$ (lenta). Para evitar divisiones, se guardan escaladas:
\begin{align}
S_F &\leftarrow S_F + p - (S_F \gg 6), & S_S &\leftarrow S_S + p - (S_S \gg 8),\\
d &= 4S_F - S_S, & \text{con } S_F = 64\,\mathrm{EMA}_r,\; S_S &= 256\,\mathrm{EMA}_l \Rightarrow d = 256(\mathrm{EMA}_r-\mathrm{EMA}_l).
\end{align}
La decisión normal usa histéresis asimétrica: $d>1280$ (diferencia mayor que 5 unidades de precio) pasa a \emph{largo}; $d<-2048$ (menor que $-8$) pasa a \emph{efectivo}; entre ambos umbrales conserva el estado. Los acumuladores se inicializan con el primer precio, sin período de arranque.

Antes de aplicar esa regla se calcula $\Delta p=p_n-p_{n-1}$. Si $\Delta p>10$, el agente pasa a largo; si $\Delta p<-10$, pasa a efectivo. Esa dirección queda bloqueada durante 1600 ticks para que el ruido inmediato no la revierta, aunque un salto grande en sentido contrario sí puede reemplazarla. Las EMA continúan actualizándose durante el bloqueo, de modo que al terminar siguen representando la historia reciente.

\textbf{Cantidad y generalidad.} En cada tick se emite comprar 63 si el estado es largo o vender 63 si está en efectivo. La maestra limita la compra a la caja disponible y la venta a las acciones poseídas; una orden cuya cantidad efectiva es cero no se ejecuta ni cobra comisión. Así se alcanza la posición posible para cualquier precio permitido entre 1 y 255, sin conocer de antemano la serie, la caja ni las acciones disponibles. El programa solo observa el flujo de precios recibido por SPI: no contiene nombres de empresas, secuencias de prueba ni decisiones precalculadas.

\subsection{Comparación con las referencias en la serie de práctica}
La Tabla~\ref{tab:refs} usa la contabilidad oficial (la misma del testbench) sobre los 20\,000 ticks de \texttt{practica.csv}; las comisiones son \$5 por orden ejecutada. La pérdida máxima se mide desde el máximo previo del patrimonio.

\begin{table}[H]\centering\small
\caption{Resultados en la serie de práctica (Meta, 20\,000 ticks).}\label{tab:refs}
\begin{tabularx}{\linewidth}{lrrrr}\toprule
Estrategia & Patrimonio final (\$) & Órdenes & Comisiones (\$) & Pérdida máx. (\$)\\\midrule
No operar & 10\,000 & 0 & 0 & 0\\
Comprar y mantener & 9\,695 & 1 & 5 & 9\,300\\
V0: SMA 128/512 & 9\,962 & 44 & 220 & 5\,163\\
\textbf{V2 adaptativa (final)} & \textbf{20\,774} & 10 & 50 & 5\,023\\\bottomrule
\end{tabularx}
\end{table}

\figopt{fig/patrimonio.pdf}{0.85}{Precio de la serie de práctica y patrimonio de cada estrategia.}

En la serie de práctica, comprar y mantener termina en \$9\,695 y V0 en \$9\,962, mientras V2 alcanza \$20\,774 con solo 10 órdenes ejecutadas. La comparación demuestra el funcionamiento sobre el conjunto público solicitado; no garantiza el mismo retorno en una serie desconocida. La robustez proviene de que las reglas dependen solo de cambios relativos y estados internos, no de valores memorizados de esta serie.

% =====================================================================
\section{Implementación en assembly}

\subsection{Estructura del programa}
Todas las variantes comparten el mismo ciclo principal: \emph{esperar precio}, \emph{actualizar estado}, \emph{predecir}, \emph{decidir} y \emph{escribir \texttt{DECISION}}. El Código~\ref{lst:asm} muestra el núcleo de la versión final \texttt{sw/trade.s}.

\begin{lstlisting}[style=asm,caption={Núcleo de \texttt{sw/trade.s} (versión final V2).},label={lst:asm}]
wait_price:                      # 1) esperar precio
    lw   x11, 0(x1)             #    STATUS del esclavo SPI
    andi x11, x11, 1            #    bit0 = precio nuevo
    beq  x11, x0, wait_price
    lw   x10, 4(x1)             #    leer PRICE (consume el aviso)
    ...                         #    (primer precio: SF = 64p, SS = 256p)
update:                          # 2) actualizar estado (EMA)
    srli x11, x5, 6
    sub  x11, x10, x11
    add  x5, x5, x11            #    SF += p - (SF >> 6)
    srli x11, x6, 8
    sub  x11, x10, x11
    add  x6, x6, x11            #    SS += p - (SS >> 8)
    slli x11, x5, 2
    sub  x11, x11, x6           # 3) d = 4*SF - SS
    sub  x12, x10, x2           # 4) cambio respecto del precio anterior
    ...                         #    salto >10 o <-10: cambia y bloquea 1600 ticks
    slti x12, x11, 1281         # 5) si no hay bloqueo, aplicar +1280/-2048
    ...
emit:                            # 6) escribir DECISION segun el estado
    addi x11, x0, 0x7F          #    comprar 63
    ...                         #    o 0xBF: vender 63
write:
    sw   x11, 8(x1)
    jal  x0, wait_price
\end{lstlisting}

\textbf{Registros (V2).} \texttt{x1}: base del periférico SPI (\texttt{0x8001\_0000}); \texttt{x2}: precio anterior; \texttt{x3}: bandera de primer precio; \texttt{x4}: ticks restantes del bloqueo; \texttt{x5}: $S_F$; \texttt{x6}: $S_S$; \texttt{x7}: estado largo/efectivo; \texttt{x10}: precio; \texttt{x11}--\texttt{x13}: temporales. El programa no usa la RAM de datos. La versión de software puro V0 usa un anillo de 512 bytes en \texttt{0x400--0x5FF} (un precio por byte, con \texttt{lbu}/\texttt{sb}) y los registros \texttt{x2} (índice), \texttt{x3} (precios vistos), \texttt{x4} ($p_0$), \texttt{x5} ($S_1$) y \texttt{x6} ($S_2$).

\textbf{Representación y operaciones faltantes.} Los precios y acumuladores son enteros de 32 bits; las diferencias $d$ y $\Delta p$ se interpretan con signo en complemento a dos. El Espino no tiene multiplicación ni división y sus shifts están deshabilitados; se resolvió así:
\begin{itemize}
  \item \textbf{V0 (sin shifts):} $128p$ y $512p$ se obtienen duplicando con \texttt{add} (9 duplicaciones, solo en el primer precio); $4S_1$ son dos \texttt{add}; las medias \emph{no se dividen}: se compara $4S_1-S_2=512(\mathrm{SMA}_{128}-\mathrm{SMA}_{512})$ con la banda de $\pm512$. Las sumas de ventana se mantienen en forma incremental con un anillo de precios.
  \item \textbf{V2 (shifter parcial):} los desplazamientos se hacen con \texttt{slli}/\texttt{srli} por 2, 6 y 8, que implementa el ALU modificado (Sección~\ref{sec:codiseno}).
\end{itemize}
El assembler del grupo se extendió con \texttt{lbu}, \texttt{sb}, \texttt{slli}, \texttt{srli} y \texttt{srai}; las codificaciones tienen pruebas unitarias (\texttt{sw/test\_assembler.py}). El programa final ocupa 49 palabras de instrucción.

\subsection{Equivalencia con la referencia}
Cada variante tiene una referencia en Python bit a bit (\texttt{ref/trader\_ref.py} para V0/V1 y \texttt{ref/trader\_ref\_ema.py} para V2) que usa el mismo modelo de retardo y la misma contabilidad del testbench. \texttt{ref/compare.py} compara tick a tick el byte devuelto por la maestra en \texttt{decisiones.csv} con la referencia:
\begin{itemize}
  \item \textbf{V2 (final):} 0 diferencias en los 20\,000 ticks de la serie de práctica.
  \item \textbf{V0:} 0 diferencias en los 20\,000 ticks.
  \item \textbf{V1:} 0 diferencias en 2000 ticks con el período de la ronda 3; en una corrida de 20\,000 ticks con período artificial de 20~$\mu$s hay una única diferencia en el tick 1, porque el acelerador rellena su anillo durante 512 ciclos con el primer precio y ese tiempo supera el presupuesto de 262 ciclos de ese período de prueba. Con los períodos reales (presupuesto $\ge 2262$ ciclos) no ocurre.
\end{itemize}
Las corridas de 20\,000 ticks usan un período de 20~$\mu$s (en lugar de 100~$\mu$s) para acortar la simulación, ya que las decisiones no dependen del período mientras el programa responda a tiempo.

\subsection{Mediciones por tick}\label{sec:medicion}
Se midió, para cada precio recibido, el número de ciclos y de instrucciones retiradas entre el fin del byte SPI (\texttt{byte\_done}) y la escritura de \texttt{DECISION}; es la misma ventana que usa el testbench oficial. El CPI es ciclos/IC en esa ventana, que incluye la espera por la detección del aviso. El reloj es de 25~MHz (40~ns por ciclo).

\begin{table}[H]\centering\small
\caption{Mediciones por tick (2000 ticks de la serie de práctica). El tiempo es ciclos $\times$ 40~ns.}\label{tab:tick}
\begin{tabularx}{\linewidth}{lrrrrrrr}\toprule
 & IC prom & IC máx & Ciclos prom & Ciclos máx & CPI prom & T prom (\textmu s) & T máx (\textmu s)\\\midrule
V0 software & 36,63 & 46 & 82,41 & 100 & 2,25 & 3,30 & 4,00\\
V1 acelerador & 21,12 & 28$^{*}$ & 51,82 & 64$^{*}$ & 2,45 & 2,07 & 2,56\\
\textbf{V2 final} & 27,51 & 33 & 61,69 & 75 & 2,24 & 2,47 & 3,00\\\bottomrule
\end{tabularx}
\par\footnotesize $^{*}$ Sin el primer tick de V1, que tarda 564 ciclos (22,6~\textmu s) por el relleno del anillo.
\end{table}

\figopt{fig/ciclos_tick.pdf}{0.7}{Distribución de los ciclos entre el precio y la decisión por tick.}

\begin{table}[H]\centering\small
\caption{Peor caso (V2) frente al período de cada ronda; presupuesto = ciclos disponibles tras la transacción SPI (9,5~\textmu s).}\label{tab:rondas}
\begin{tabularx}{\linewidth}{lrrrr}\toprule
Ronda & Período $T$ & Presupuesto (ciclos) & Peor caso / $T$ & Peor caso / presupuesto\\\midrule
1 & 10 ms & 249\,762 & 0,030\,\% & 0,030\,\%\\
2 & 1 ms & 24\,762 & 0,30\,\% & 0,30\,\%\\
3 & 0,1 ms & 2\,262 & 3,0\,\% & 3,3\,\%\\\bottomrule
\end{tabularx}
\end{table}
El peor caso de todas las variantes está muy por debajo del presupuesto de la ronda 3, de modo que el programa siempre alcanza a escribir la decisión antes del siguiente tick: no se repite nunca la decisión anterior por falta de tiempo y no se pierde ningún precio (\texttt{precios\_perdidos = 0}). El testbench oficial pasa sus 4 chequeos con la versión final.

% =====================================================================
\section{Co-diseño hardware/software}\label{sec:codiseno}

\subsection{Qué consume más ciclos: perfil y ley de Amdahl}
Se perfiló la variante V0 con un contador de ciclos por instrucción (\texttt{tb/tb\_metrics.sv} escribe \texttt{perfil\_pc.csv}). De los 82,4 ciclos por tick, el tramo que \emph{actualiza el anillo de precios y calcula $d$} (desde \texttt{update} hasta $d=4S_1-S_2$) consume el 56,3\,\%; la espera del aviso, la histéresis y el armado de la decisión se reparten el resto. Por la ley de Amdahl, acelerar ese tramo infinitamente rápido da un speedup máximo $1/(1-0{,}563)=2{,}29\times$.

Un speedup real depende de lo que cuesta el reemplazo: en V1 el procesador debe escribir el precio al periférico, esperar que termine (5 ciclos) y leer $d$; en V2 el tramo no desaparece sino que se reduce a 11 instrucciones ALU gracias a que la estrategia EMA no necesita anillo.

\begin{table}[H]\centering\small
\caption{Speedup frente a la versión de software puro V0.}\label{tab:amdahl}
\begin{tabularx}{\linewidth}{lrrY}\toprule
Variante & Predicho (cota Amdahl) & Medido & Comentario\\\midrule
V1 acelerador & $\le 2{,}29\times$ & $1{,}59\times$ & Queda el sondeo del acelerador (15 ciclos/tick) y la transferencia por MMIO.\\
V2 shifter + EMA & $\le 2{,}29\times$ & $1{,}34\times$ & La detección de saltos agrega control, pero evita el anillo de precios.\\\bottomrule
\end{tabularx}
\end{table}
El speedup medido queda bajo la cota porque la parte no acelerable (esperar el aviso del SPI y escribir \texttt{DECISION}) pesa cerca del 44\,\% del tiempo y porque el reemplazo no es gratis. El CPI sube en V1 (2,45) porque quedan relativamente más cargas MMIO, que toman 2 ciclos.

\subsection{Variantes implementadas}
\begin{itemize}
  \item \textbf{V0 (base, solo software).} \texttt{sw/trade\_sw.s}: SMA 128/512 con anillo de 512 bytes en RAM. No cambia el RTL.
  \item \textbf{V1 (periférico MMIO).} \texttt{rtl/pochoco\_sma.v}, en \texttt{0x8002\_0000}: guarda el anillo de 512 precios en un bloque de RAM y mantiene $d=4S_1-S_2$ de forma incremental, $d\leftarrow d+3p+p_{-512}-4p_{-128}$, con un único sumador/restador de 18 bits en 5 ciclos. Con el primer precio rellena el anillo con $p_0$ (512 ciclos). El Espino lee $d$, aplica la histéresis y escribe \texttt{DECISION}: el hardware solo acelera una parte y no reemplaza al procesador. Tiene testbench propio (\texttt{tb/tb\_sma.sv}), que compara $d$ con el cálculo directo en 3000 ticks (0 errores). Un diseño inicial con dos sumas independientes ocupaba 1456 celdas lógicas y no cabía; la versión compacta usa un sumador compartido y quita el contador de precios.
  \item \textbf{V2 (extensión del ALU).} Parámetro \texttt{UseShifter} del \texttt{espino\_alu}: ejecuta \texttt{SLL} por 2, 6 y 8 y \texttt{SRL} por 6 y 8; cualquier otro desplazamiento (y \texttt{SRA}) sigue ejecutándose como \texttt{ADD}, igual que el Espino base. Un barrel shifter completo de 32 bits costaba $\approx$340 LUTs y no cabía; el parcial cuesta $\approx$180. Se acompaña de la estrategia EMA, que sin shifter no se puede ejecutar. Se verificó con el testbench oficial y la comparación contra la referencia.
\end{itemize}
Las variantes V1 y V2 usan además \texttt{LitePeriph=1}, que elimina LEDs, displays y el contador de ciclos del SoC (no se usan en el trader) y libera $\approx$160 celdas lógicas; sin esa reducción ninguna de las dos cabe en las 1280 celdas del iCE40 HX1K. La Tabla~\ref{tab:hw} incluye V0b (V0 con esos periféricos reducidos) para comparar las variantes en igualdad de condiciones.

\begin{table}[H]\centering\small
\caption{Variantes: recursos de síntesis (Yosys + nextpnr, \texttt{--hx1k --package vq100 --freq 25}) y desempeño medido.}\label{tab:hw}
\begin{tabularx}{\linewidth}{lrrrrrrr}\toprule
Variante & LUTs & Celdas lóg. (1280) & BRAM (16) & Fmax (MHz) & Ciclos prom/máx & IC prom & CPI\\\midrule
V0 software & 973 & 1152 & 12 & 46,8 & 82,4 / 100 & 36,6 & 2,25\\
V0b software, periféricos lite & 868 & 1033 & 12 & 47,0 & 82,4 / 100 & 36,6 & 2,25\\
V1 acelerador SMA & 994 & 1195 & 13 & 42,0 & 51,8 / 564$^{*}$ & 21,1 & 2,45\\
\textbf{V2 shifter + EMA} & 1040 & 1205 & 12 & 45,3 & 61,7 / 75 & 27,5 & 2,24\\\bottomrule
\end{tabularx}
\par\footnotesize $^{*}$ Primer tick (relleno del anillo); el resto, máx 64.
\end{table}
Las tres variantes cierran el timing a 25~MHz con margen (Fmax $\ge 42$~MHz). Una cuarta alternativa, subir la frecuencia con la PLL, \textbf{no es posible en esta placa}: el iCE40 HX1K de la Go Board no tiene PLL y nextpnr rechaza la primitiva (\emph{no suitable BEL}). Por lo tanto el presupuesto de ciclos de cada ronda se mantiene en el de 25~MHz (Tabla~\ref{tab:rondas}).

\subsection{Ops/s y roofline}
Se define \textbf{operación útil} como una instrucción aritmética o lógica del kernel (\texttt{add}, \texttt{sub}, \texttt{slli}, \texttt{srli}, \texttt{slti}, etc.) ejecutada en la ventana de cálculo; se excluye el sondeo del aviso (la espera). En V1 se suman las 4 sumas del acelerador. El \textbf{tráfico} es la suma de bytes de las cargas y almacenamientos del kernel (\texttt{lw}/\texttt{sw} 4~B, \texttt{lbu}/\texttt{sb} 1~B). Con ello:

\begin{table}[H]\centering\small
\caption{Rendimiento e intensidad operacional por tick (promedio).}\label{tab:roof}
\begin{tabularx}{\linewidth}{lrrrr}\toprule
Variante & Ops útiles/tick & Bytes/tick & Intensidad (ops/B) & Ops/s (Mops/s)\\\midrule
V0 & 19,5 & 10,7 & 1,83 & 6,41\\
V1 & 7,4 & 16,0 & 0,46 & 5,14\\
\textbf{V2} & 15,2 & 8,0 & 1,90 & 6,77\\\bottomrule
\end{tabularx}
\end{table}
Ops/s es $\text{ops/tick}$ dividido por el tiempo de cálculo del tick (ciclos de la ventana menos los de sondeo, a 25~MHz). Es una medida de qué tan bien se usa el procesador durante el cálculo; no es comparable entre variantes que hacen un número distinto de operaciones por tick, por lo que la métrica para comparar variantes es el tiempo por tick de la Tabla~\ref{tab:tick}. Una variante con menos operaciones (V2) puede resolver el mismo problema más rápido con un Ops/s parecido.

El techo del Espino se modela con $P_{max}=f\times 1\ \text{op/ciclo}=25$~Mops/s (CPI ideal 1) y un ancho de banda de 2~B/ciclo (una carga ocupa 2 ciclos), es decir 50~MB/s; el punto de quiebre está en $0{,}5$~ops/B.

\figopt{fig/roofline.pdf}{0.7}{Roofline del Espino (25~MHz) con los kernels de V0, V1 y V2.}

V0 y V2 están en la zona limitada por cómputo, con 6,41 y 6,77~Mops/s: el CPI real cercano a 2,24 y los saltos y cargas con latencia impiden acercarse al techo de 25~Mops/s. V1 queda en la zona limitada por el bus (intensidad de 0,46~ops/B, bajo el punto de quiebre de 0,5): el kernel dejó de hacer aritmética y quedó dominado por las transferencias MMIO al acelerador.

\subsection{Elección de la variante final}
Se eligió \textbf{V2 (shifter parcial + EMA)}:
\begin{itemize}
  \item \textbf{Tiempo ampliamente suficiente y sin latencia de arranque.} Promedia 61,7 ciclos y su peor caso es 75 ciclos; V1 es más rápida en régimen, pero tarda 564 ciclos en el primer tick. Ambas cumplen holgadamente la ronda 3.
  \item \textbf{Mejor resultado algorítmico en la práctica.} Alcanza \$20\,774 contra \$9\,962 de V0, con 10 órdenes ejecutadas en vez de 44. La detección de saltos y la histéresis reducen cambios innecesarios.
  \item \textbf{Equilibrio de área.} 1040 LUTs y 1205 celdas lógicas (94\,\%), 12 de 16 BRAM, sin anillo de precios en RAM y con Fmax de 45~MHz. V1 consume un BRAM más y tiene la Fmax más baja (42~MHz).
\end{itemize}
No es la de menor área (V0b) ni la de menor promedio de ciclos (V1), pero equilibra recursos, tiempo de respuesta y calidad de decisión, y es la única en que el cambio de hardware habilita directamente la estrategia final. Como todas responden dentro del presupuesto, la elección se decide principalmente por el resultado del agente y la ausencia de un arranque costoso.

% =====================================================================
\section{Verificación y reproducibilidad}\label{sec:repro}
Verificación realizada: (i) pruebas unitarias del assembler (8/8); (ii) testbench oficial \texttt{tb\_proyecto\_4.sv} con la versión final (4/4); (iii) \texttt{tb\_sma.sv} para el acelerador; (iv) comparación bit a bit contra la referencia en Python, con 0 diferencias en 20\,000 ticks; (v) síntesis y place \& route a 25~MHz de las cuatro configuraciones, sin errores de timing.

\begin{lstlisting}[style=shell]
make trade                         # ensambla trade.s, trade_sw.s, trade_hw.s con sw/assembler.py
make p4-test                       # testbench oficial con la version final
make p4-full-check                 # 20 000 ticks y comparacion exacta: 0 diferencias
bash ref/hw_report.sh              # tabla de variantes (LUTs, Fmax, ciclos, IC, CPI)
make all                           # sintetiza, P&R a 25 MHz, bitstream y programa la Go Board
\end{lstlisting}
Herramientas: Yosys 0.68+130, nextpnr-ice40 0.11.1, Icarus Verilog 14.0, Python 3.9.6.

% =====================================================================
\section{Conclusiones}
\begin{itemize}
  \item El diseño final combina dos EMA, histéresis asimétrica y detección de saltos. No usa multiplicación ni división, no contiene datos conocidos de las pruebas y opera sobre cualquier precio permitido por el protocolo.
  \item La implementación coincide exactamente con su referencia Python en los 20\,000 ticks. En la serie pública alcanza \$20\,774, frente a \$10\,000 sin operar, \$9\,695 al comprar y mantener y \$9\,962 de la versión V0.
  \item El peor caso es 75 ciclos (3,00~$\mu$s), equivalente al 3,3\,\% del presupuesto disponible de la ronda más exigente. Por lo tanto existe margen suficiente para recibir cada precio y actualizar la decisión a tiempo.
  \item El tramo acelerable de V0 representa el 56,3\,\% de los ciclos y fija una cota de Amdahl de $2{,}29\times$. V1 logra $1{,}59\times$ y V2 $1{,}34\times$; V2 agrega control para la estrategia adaptativa, pero mantiene Fmax de 45,3~MHz y cabe en el iCE40 HX1K.
  \item La serie de práctica permite verificar y comparar, pero no garantiza el retorno en datos desconocidos. La justificación de la alternativa final combina corrección, generalidad, recursos, tiempo de respuesta y resultado observado, sin afirmar que una serie futura tendrá el mismo comportamiento.
\end{itemize}
\end{document}
